\documentclass[10pt,a4paper]{article}
\usepackage[margin=0.72in,headheight=15pt]{geometry}
\usepackage[T1]{fontenc}
\usepackage{lmodern}
\usepackage{xcolor}
\usepackage{listings}
\usepackage{enumitem}
\usepackage{booktabs,tabularx,longtable}
\usepackage{fancyhdr}
\usepackage{hyperref}
\usepackage{microtype}
\usepackage{amssymb}
\usepackage[most]{tcolorbox}
\hypersetup{colorlinks=true,linkcolor=blue!55!black,urlcolor=blue!55!black}
\definecolor{navy}{HTML}{16324F}
\definecolor{codebg}{HTML}{F5F7FA}
\pagestyle{fancy}\fancyhf{}\lhead{CSE341 -- 8086 Lab Mid}\rhead{Solved Exercises / Code Templates}\cfoot{\thepage}
\setlist{nosep,leftmargin=*}\setlength{\parskip}{3pt}\setlength{\parindent}{0pt}
\lstdefinestyle{asm}{language={},basicstyle=\ttfamily\footnotesize,backgroundcolor=\color{codebg},frame=single,rulecolor=\color{black!15},breaklines=true,columns=fullflexible,keepspaces=true,showstringspaces=false,tabsize=4}
\newtcolorbox{noteBox}{colback=yellow!10,colframe=orange!60!black,boxrule=.6pt,arc=2mm}
\newtcolorbox{answerBox}{colback=green!5,colframe=green!40!black,boxrule=.6pt,arc=2mm}
\title{\Huge\bfseries 8086 Assembly Language\\[3pt]\Large Lab Midterm Solutions and Template Workbook\\[5pt]\normalsize Lab 01 through Lab 03(B) + Sample Midterm 1 and 2}
\author{CSE341 -- Microprocessors}
\date{Companion to the 8086 preparation workbook}
\begin{document}
\maketitle
\begin{tcolorbox}[colback=navy!5,colframe=navy,title=How to use these answers]
This file follows the supplied \texttt{Assembly Template.docx}. Unless a section says ``complete program'', code is organized as \textbf{DATA declarations} plus a \textbf{MAIN body}. Put declarations inside \texttt{.DATA}; put the main body after DS initialization and before the standard exit sequence. If a helper procedure is shown, place it after \texttt{MAIN ENDP} and before \texttt{END MAIN}. Each exercise is independent: do not paste every answer into one program.

Most solutions use MASM/EMU8086-style syntax and DOS \texttt{INT 21H}. Values loaded directly into registers stand in for the stated input/initial-register condition. Where the lab statement is ambiguous or mathematically inconsistent, the assumption/correction is identified.
\end{tcolorbox}
\tableofcontents
\newpage

\section{Common template and reusable procedures}
\subsection{Base template from the supplied DOCX}
\begin{lstlisting}[style=asm]
.MODEL SMALL
.STACK 100H
.DATA
    ; declare variables and strings here
.CODE
MAIN PROC
    MOV AX, @DATA
    MOV DS, AX

    ; paste one exercise's MAIN BODY here

    MOV AX, 4C00H
    INT 21H
MAIN ENDP
END MAIN
\end{lstlisting}

\subsection{Print an unsigned 16-bit decimal number}
This procedure expects the number in AX. It handles 0 and prints digits in normal order by pushing remainders and popping them. It preserves BX, CX, DX.
\begin{lstlisting}[style=asm]
PRINT_U16 PROC
    PUSH BX
    PUSH CX
    PUSH DX
    CMP AX, 0
    JNE PU_CONVERT
    MOV DL, '0'
    MOV AH, 2
    INT 21H
    JMP PU_DONE
PU_CONVERT:
    XOR CX, CX
    MOV BX, 10
PU_DIVIDE:
    XOR DX, DX
    DIV BX              ; AX=quotient, DX=remainder
    PUSH DX
    INC CX
    CMP AX, 0
    JNE PU_DIVIDE
PU_OUTPUT:
    POP DX
    ADD DL, '0'
    MOV AH, 2
    INT 21H
    LOOP PU_OUTPUT
PU_DONE:
    POP DX
    POP CX
    POP BX
    RET
PRINT_U16 ENDP
\end{lstlisting}

\subsection{Print a signed 16-bit decimal number}
Expects the signed value in AX. The magnitude of -32768 is 32768, which the unsigned helper can print.
\begin{lstlisting}[style=asm]
PRINT_S16 PROC
    PUSH AX
    TEST AX, 8000H
    JZ PS_POSITIVE
    MOV DL, '-'
    MOV AH, 2
    INT 21H
    POP AX
    NEG AX
    CALL PRINT_U16
    RET
PS_POSITIVE:
    POP AX
    CALL PRINT_U16
    RET
PRINT_S16 ENDP
\end{lstlisting}

\subsection{Print CR/LF}
\begin{lstlisting}[style=asm]
NEWLINE PROC
    MOV DL, 0DH
    MOV AH, 2
    INT 21H
    MOV DL, 0AH
    MOV AH, 2
    INT 21H
    RET
NEWLINE ENDP
\end{lstlisting}

\begin{noteBox}
When a program uses a procedure, make sure the procedure is placed after \texttt{MAIN ENDP}. The procedures above are optional helpers; the lab's simplest one-character tasks do not need them. All code blocks are standalone exercise answers or bodies, not one single combined program.
\end{noteBox}

\section{Lab 01 -- Introduction and basic arithmetic}
\subsection*{Task 01 -- Move AX to BX}
\textbf{Given:} value already in AX. The problem is a register transfer, not keyboard text input.
\begin{lstlisting}[style=asm]
; MAIN body
MOV AX, 1234H       ; sample initial value
MOV BX, AX          ; BX receives AX
\end{lstlisting}
\textit{If the examiner literally asks for a keyboard character first:} use \texttt{AH=01H}, then zero-extend the ASCII byte if treating it as a byte value: \texttt{MOV AH,0 / MOV BX,AX}. Convert from ASCII first if it is a decimal digit.

\subsection*{Task 02 -- Swap two numbers using at most three registers}
\begin{lstlisting}[style=asm]
MOV AX, 10
MOV BX, 20
MOV CX, AX
MOV AX, BX
MOV BX, CX
; AX=20, BX=10
\end{lstlisting}

\subsection*{Task 03 -- Add two numbers using two registers}
\begin{lstlisting}[style=asm]
MOV AX, 25
MOV BX, 17
ADD AX, BX          ; AX=42
\end{lstlisting}

\subsection*{Task 04 -- Subtract two numbers}
\begin{lstlisting}[style=asm]
MOV AX, 3
MOV BX, 5
SUB AX, BX          ; AX=FFFEH (signed -2)
; For unsigned subtraction, CF=1 indicates a borrow.
\end{lstlisting}
The numeric bit pattern is fixed-width. A negative signed answer appears as its two's-complement representation. It is not an ordinary unsigned magnitude.

\subsection*{Task 05 -- Swap using ADD/SUB only}
\begin{lstlisting}[style=asm]
MOV AX, 5
MOV BX, 10
ADD AX, BX          ; AX = old AX + old BX
SUB BX, AX          ; BX = -old AX
NEG BX              ; BX = old AX
SUB AX, BX          ; AX = old BX
; AX=10, BX=5
\end{lstlisting}
This method can overflow if the sum cannot fit in the register width. A temporary register or XCHG is safer in ordinary programming.

\subsection*{Task 06 -- Arithmetic expressions with variables}
\begin{lstlisting}[style=asm]
.DATA
A DW 10
B DW 5
C DW ?
.CODE
; (1) A = B - A
MOV AX, B
SUB AX, A
MOV A, AX

; (2) A = -(A + 1)
MOV AX, A
INC AX
NEG AX
MOV A, AX

; (3) C = A + (B + 1), using INC
MOV AX, B
INC AX
ADD AX, A
MOV C, AX

; (4) A = B - (A - 1), using DEC
MOV AX, A
DEC AX
MOV BX, B
SUB BX, AX
MOV A, BX
\end{lstlisting}
The expressions run sequentially, so each later part uses the updated A from the earlier part. Reset initial A and B before testing a specific sub-expression.

\subsection*{Task 07 -- Multiplication and division}
\begin{lstlisting}[style=asm]
.DATA
X DW 6
Y DW 3
Z DW 2
RESULT DW ?
.CODE
; (1) X * Y
MOV AX, X
MUL Y               ; DX:AX = X*Y
MOV RESULT, AX      ; low word if result fits in 16 bits

; (2) X / Y
MOV AX, X
XOR DX, DX
DIV Y               ; AX=quotient, DX=remainder
MOV RESULT, AX

; (3) X * Y / Z
MOV AX, X
MUL Y               ; preserve full product in DX:AX
DIV Z               ; divide 32-bit product by Z
MOV RESULT, AX      ; quotient
\end{lstlisting}
\textbf{Assumption:} Z is nonzero and final quotient fits in AX. The multiplication result is 32-bit; do not clear DX between MUL and DIV in part (3).

\subsection*{Task 08 -- Hexadecimal arithmetic}
\begin{enumerate}[label=\alph*)]
\item \texttt{36DFH * AFH}: load AX=36DFH and BX=00AFH. Full product is DX:AX = \texttt{0025:8271H}.
\item \texttt{F4D5H / C9A5H}: load AX=F4D5H, DX=0, BX=C9A5H; quotient AX=0001H and remainder DX=2B30H.
\item \texttt{CA92H * BAF9H}: full product is DX:AX = \texttt{93F3:1C02H}.
\item \texttt{C2A2H * ABCDH / BEDH}: the quotient is 000AF3D8H (20 bits) and remainder 0AC2H. A single 16-bit DIV cannot hold this quotient in AX; use the 32-bit-by-16-bit long-division routine below.
\end{enumerate}
\begin{lstlisting}[style=asm]
; Reusable instruction pattern for each expression
MOV AX, 36DFH
MOV BX, 00AFH
MUL BX              ; DX:AX = 36DFH * 00AFH

MOV AX, 0F4D5H
XOR DX, DX
MOV BX, 0C9A5H
DIV BX              ; AX=1, DX=2B30H

MOV AX, 0CA92H
MOV BX, 0BAF9H
MUL BX              ; full result DX:AX

MOV AX, 0C2A2H
MOV BX, 0ABCDH
MUL BX
MOV PRODHI,DX        ; high word of product
MOV PRODLO,AX        ; low word of product
MOV DIVISOR,0BEDH
; 32-bit dividend / 16-bit divisor, quotient may need 32 bits
MOV QHI,0
MOV QLO,0
MOV REM16,0
MOV CX,32
T8_DIV32:
    SHL PRODLO,1
    RCL PRODHI,1     ; CF = next dividend bit
    RCL REM16,1      ; REM16 = REM16*2 + next bit
    SHL QLO,1
    RCL QHI,1        ; shift quotient left
    MOV AX,REM16
    CMP AX,DIVISOR
    JB T8_NO_SUB
    SUB AX,DIVISOR
    MOV REM16,AX
    OR QLO,1
T8_NO_SUB:
    LOOP T8_DIV32
; Expected QHI:QLO = 000A:F3D8H, REM16 = 0AC2H
\end{lstlisting}
\textbf{Add these declarations for the last subproblem:} \texttt{PRODHI, PRODLO, QHI, QLO, REM16, DIVISOR} as WORD variables. This long division works here because the intermediate remainder stays below 2*0BEDH and fits in 16 bits. Do not execute the ordinary \texttt{DIV CX} for this subproblem: the quotient 0AF3D8H exceeds FFFFH and would cause divide overflow.

\subsection*{Task 09 -- Two MOV examples for each common operand combination}
Assume these variables are declared: \texttt{WORD1 DW ?} and \texttt{BYTE1 DB ?}.
\begin{lstlisting}[style=asm]
; Register <- register
MOV AX, BX
MOV DL, CL
; Register <- memory
MOV AX, WORD1
MOV BL, BYTE1
; Memory <- register
MOV WORD1, AX
MOV BYTE1, BL
; Register <- immediate
MOV CX, 1234H
MOV AL, 25
; Memory <- immediate
MOV WORD1, 100
MOV BYTE1, 7
\end{lstlisting}
No memory-to-memory MOV is allowed. Also avoid mixing byte and word sizes.

\subsection*{Task 10 -- Two ADD and SUB examples for each common combination}
Use the same \texttt{WORD1} and \texttt{BYTE1} declarations from Task 09.
\begin{lstlisting}[style=asm]
; Register, register
ADD AX, BX
SUB CL, DL
; Register, memory
ADD AX, WORD1
SUB AL, BYTE1
; Memory, register
ADD WORD1, BX
SUB BYTE1, CL
; Register, immediate
ADD AX, 10
SUB DL, 2
; Memory, immediate
ADD WORD1, 5
SUB BYTE1, 1
\end{lstlisting}
The destination is always updated. Memory-to-memory arithmetic is invalid; move one memory operand into a register first.

\subsection*{Task 11 -- Evaluate (1+2)*(3-1)/5+3+2-(1*2)}
Integer division is used. The expression evaluates as $3\times2=6$, $6/5=1$, then $1+3+2-2=4$.
\begin{lstlisting}[style=asm]
MOV AX, 1
ADD AX, 2           ; 3
MOV BX, AX
MOV AX, 3
SUB AX, 1           ; 2
MUL BX              ; DX:AX = 6
MOV BX, 5
DIV BX              ; AX = 1
ADD AX, 3           ; 4
ADD AX, 2           ; 6
MOV BX, 1
MOV CX, 2
MOV AX, BX
MUL CX              ; AX = 2
MOV BX, AX
; Restore the partial expression result if it was saved:
; preferred complete approach uses stack:
\end{lstlisting}
\begin{lstlisting}[style=asm]
; Complete version with PUSH/POP preserving the partial result
MOV AX, 1
ADD AX, 2           ; 3
MOV BX, AX
MOV AX, 3
SUB AX, 1           ; 2
MUL BX              ; product = 6 in AX (fits)
MOV BX, 5
XOR DX, DX
DIV BX              ; integer quotient = 1
ADD AX, 3
ADD AX, 2           ; partial = 6
PUSH AX
MOV AX, 1
MOV BX, 2
MUL BX              ; 1*2 = 2
MOV BX, AX
POP AX
SUB AX, BX          ; final AX=4
\end{lstlisting}

\section{Lab 02 -- Basic I/O, arithmetic, and flags}
\subsection*{Task 01 -- Prompt, read, and display one character}
\begin{lstlisting}[style=asm]
.DATA
MSG DB 'Please insert a character: $'
.CODE
; MAIN body
LEA DX, MSG
MOV AH, 9
INT 21H
MOV AH, 1
INT 21H
MOV BL, AL          ; preserve input
MOV DL, BL
MOV AH, 2
INT 21H
\end{lstlisting}

\subsection*{Task 02 -- Input two digits; show add, subtract, multiply, divide}
This complete MAIN body accepts one decimal digit per operand. Subtraction is displayed as a signed difference; division shows both quotient and remainder. Add the shown data declarations and use the shared \texttt{PRINT\_U16} helper.
\begin{lstlisting}[style=asm]
.DATA
P1 DB 'First digit (0-9): $'
P2 DB 0DH,0AH,'Second digit (0-9): $'
M1 DB 0DH,0AH,'ADD = $'
M2 DB 0DH,0AH,'SUB = $'
M3 DB 0DH,0AH,'MUL = $'
M4 DB 0DH,0AH,'DIV quotient = $'
M5 DB ', remainder = $'
ERR DB 0DH,0AH,'Cannot divide by zero.$'
A DB ?
B DB ?
Q DB ?
REMV DB ?
.CODE
LEA DX,P1
MOV AH,9
INT 21H
MOV AH,1
INT 21H
SUB AL,'0'
MOV A,AL
LEA DX,P2
MOV AH,9
INT 21H
MOV AH,1
INT 21H
SUB AL,'0'
MOV B,AL

LEA DX,M1
MOV AH,9
INT 21H
XOR AX,AX
MOV AL,A
XOR BX,BX
MOV BL,B
ADD AX,BX
CALL PRINT_U16

LEA DX,M2
MOV AH,9
INT 21H
MOV AL,A
CMP AL,B
JAE T2_SUB_POS
MOV DL,'-'
MOV AH,2
INT 21H
MOV AL,B
SUB AL,A
JMP T2_SUB_PRINT
T2_SUB_POS:
MOV AL,A
SUB AL,B
T2_SUB_PRINT:
XOR AH,AH
CALL PRINT_U16

LEA DX,M3
MOV AH,9
INT 21H
MOV AL,A
MOV BL,B
MUL BL
CALL PRINT_U16

CMP B,0
JE T2_DIV_ZERO
LEA DX,M4
MOV AH,9
INT 21H
MOV AL,A
XOR AH,AH
MOV BL,B
DIV BL              ; AL=quotient, AH=remainder
MOV Q,AL
MOV REMV,AH         ; preserve remainder before DOS calls
XOR AX,AX
MOV AL,Q
CALL PRINT_U16
LEA DX,M5
MOV AH,9
INT 21H
XOR AX,AX
MOV AL,REMV
CALL PRINT_U16
JMP T2_FINISH
T2_DIV_ZERO:
LEA DX,ERR
MOV AH,9
INT 21H
T2_FINISH:
\end{lstlisting}
\subsection*{Task 03(a) -- Read a character and display it at the next same-line position}
\begin{lstlisting}[style=asm]
MOV AH, 1
INT 21H
MOV BL, AL
MOV DL, BL
MOV AH, 2
INT 21H
\end{lstlisting}
\subsection*{Task 03(b) -- Read uppercase and display lowercase on the same line}
\begin{lstlisting}[style=asm]
MOV AH, 1
INT 21H
ADD AL, 20H         ; assumes input is A..Z
MOV DL, AL
MOV AH, 2
INT 21H
\end{lstlisting}
\subsection*{Task 04 -- Uppercase to lowercase on the next line}
\begin{lstlisting}[style=asm]
MOV AH, 1
INT 21H
ADD AL, 20H
MOV BL, AL
MOV DL, 0DH
MOV AH, 2
INT 21H
MOV DL, 0AH
INT 21H
MOV DL, BL
INT 21H
\end{lstlisting}

\subsection*{Task 05 -- Add two digits and show the message}
Assumption from the sheet: both inputs are single decimal digits and their sum is below 10.
\begin{lstlisting}[style=asm]
.DATA
MSG1 DB 0DH,0AH,'THE SUM OF $'
MSG2 DB ' AND $'
MSG3 DB ' IS $'
.CODE
MOV DL,'?'
MOV AH,2
INT 21H
MOV AH,1
INT 21H
MOV BL,AL          ; original ASCII digit
SUB AL,'0'
MOV CL,AL
MOV AH,1
INT 21H
MOV BH,AL          ; second ASCII digit
SUB AL,'0'
ADD CL,AL          ; numeric sum
LEA DX,MSG1
MOV AH,9
INT 21H
MOV DL,BL
MOV AH,2
INT 21H
LEA DX,MSG2
MOV AH,9
INT 21H
MOV DL,BH
MOV AH,2
INT 21H
LEA DX,MSG3
MOV AH,9
INT 21H
MOV DL,CL
ADD DL,'0'
MOV AH,2
INT 21H
\end{lstlisting}

\subsection*{Task 06 -- Display three initials vertically}
\begin{lstlisting}[style=asm]
.DATA
MSG DB 'ENTER THREE INITIALS: $'
.CODE
LEA DX,MSG
MOV AH,9
INT 21H
MOV AH,1
INT 21H
MOV BL,AL
MOV AH,1
INT 21H
MOV BH,AL
MOV AH,1
INT 21H
MOV CL,AL
; first initial
CALL NEWLINE
MOV DL,BL
MOV AH,2
INT 21H
; second initial
CALL NEWLINE
MOV DL,BH
INT 21H
; third initial
CALL NEWLINE
MOV DL,CL
INT 21H
\end{lstlisting}

\subsection*{Task 07 -- Convert hex digit A--F to decimal}
\begin{lstlisting}[style=asm]
.DATA
P DB 'ENTER A HEX DIGIT (A-F): $'
O DB 0DH,0AH,'IN DECIMAL IT IS $'
.CODE
LEA DX,P
MOV AH,9
INT 21H
MOV AH,1
INT 21H
SUB AL,37H          ; 'A'=41H -> 10, ..., 'F' -> 15
XOR AH,AH
MOV BL,AL
LEA DX,O
MOV AH,9
INT 21H
XOR AX,AX
MOV AL,BL
CALL PRINT_U16
\end{lstlisting}
This task as written accepts A--F. The later Lab 03(B) Task 06 expands validation to accept 0--9 and A--F.

\subsection*{Task 08 -- Display a solid 10 x 10 box of asterisks}
\begin{lstlisting}[style=asm]
.DATA
LINE DB '**********',0DH,0AH,'$'
.CODE
MOV CX,10
PRINT_ROW:
    LEA DX,LINE
    MOV AH,9
    INT 21H
    LOOP PRINT_ROW
\end{lstlisting}

\subsection*{Task 09 -- Read two multi-digit decimal numbers and add them}
Use a helper that reads decimal digits until Enter (CR). The result is a 16-bit unsigned sum, so each input and the sum must fit 0..65535. Insert the data declarations and main body below; put \texttt{READ\_DECIMAL} and \texttt{PRINT\_U16} procedures after MAIN ENDP.
\begin{lstlisting}[style=asm]
.DATA
PROMPT1 DB 'First number: $'
PROMPT2 DB 0DH,0AH,'Second number: $'
OUTMSG DB 0DH,0AH,'Sum = $'
N1 DW ?
N2 DW ?
SUM DW ?
.CODE
LEA DX,PROMPT1
MOV AH,9
INT 21H
CALL READ_DECIMAL
MOV N1,AX
LEA DX,PROMPT2
MOV AH,9
INT 21H
CALL READ_DECIMAL
MOV N2,AX
MOV AX,N1
ADD AX,N2
MOV SUM,AX
LEA DX,OUTMSG
MOV AH,9
INT 21H
MOV AX,SUM
CALL PRINT_U16
\end{lstlisting}
\begin{lstlisting}[style=asm]
READ_DECIMAL PROC
    PUSH BX
    PUSH CX
    PUSH DX
    XOR BX,BX            ; accumulator
RD_NEXT:
    MOV AH,1
    INT 21H
    CMP AL,0DH
    JE RD_DONE
    CMP AL,'0'
    JB RD_NEXT           ; ignore non-digit input
    CMP AL,'9'
    JA RD_NEXT
    SUB AL,'0'
    XOR AH,AH
    MOV CX,AX            ; keep current digit
    MOV AX,BX
    MOV DX,10
    MUL DX               ; accumulator * 10
    ADD AX,CX
    MOV BX,AX
    JMP RD_NEXT
RD_DONE:
    MOV AX,BX
    POP DX
    POP CX
    POP BX
    RET
READ_DECIMAL ENDP
\end{lstlisting}

\subsection*{Task 10 -- Flags after AL=80H, BL=80H, ADD AL,BL}
\begin{answerBox}
Result stored in AL: \texttt{00H}. Flags: \textbf{CF=1, PF=1, AF=0, ZF=1, SF=0, OF=1}. Explanation: 80H+80H=100H; the 8-bit carry is discarded, giving zero. Two negative signed operands (-128 and -128) produce an out-of-range mathematical result, so OF=1. The low nibbles 0+0 create no auxiliary carry.
\end{answerBox}
The notes' worked line lists AF=1, but that is a mistake: AF is 0 for this addition.

\subsection*{Task 11 -- Both AX and BX positive: carry in but no carry out iff signed overflow}
For positive signed 16-bit numbers, both sign bits are 0. If their sum crosses the positive maximum (7FFFH), the 16-bit result's sign bit becomes 1. This requires a carry into bit 15 but no carry out of bit 15. Since OF is the XOR of carry into and carry out of the sign bit, OF=1 exactly in this case. Conversely, if the carries differ while both inputs have sign 0, the carry pattern must be carry-in=1 and carry-out=0, so the stored result has sign bit 1 and signed overflow has occurred.

\subsection*{Task 12 -- Both AX and BX negative: carry out but no carry in iff signed overflow}
For two negative signed 16-bit numbers, both sign bits are 1. If the sum is below -32768, the stored 16-bit result becomes non-negative. The carry pattern is carry-in=0 and carry-out=1. Again, OF=carry-in XOR carry-out, so overflow occurs exactly when carry-out exists but carry-in does not. Conversely, with both input signs 1, that unequal carry pattern is the condition that makes the result sign bit 0 and indicates signed overflow.

\section{Lab 03(A) -- Flow control and branching}
Unless otherwise stated, the snippets are the MAIN body; define the named variables/messages in .DATA. Use signed comparisons for signed values and unsigned comparisons for non-negative values/ASCII codes.

\subsection*{Task 01 -- Replace AX by 5 if AX is negative}
\begin{lstlisting}[style=asm]
MOV AX, -10          ; sample input
CMP AX, 0
JGE T1_DONE
MOV AX, 5
T1_DONE:
\end{lstlisting}
\subsection*{Task 02 -- Display the earlier extended ASCII character in AL/BL}
Use unsigned ordering because character codes are byte values.
\begin{lstlisting}[style=asm]
MOV AL, 'B'
MOV BL, 'A'
CMP AL, BL
JBE T2_AL_FIRST
MOV DL, BL
JMP T2_PRINT
T2_AL_FIRST:
MOV DL, AL
T2_PRINT:
MOV AH, 2
INT 21H
\end{lstlisting}
\subsection*{Task 03 -- Put -1, 0, or 1 in BX according to AX}
\begin{lstlisting}[style=asm]
CMP AX, 0
JL T3_NEG
JE T3_ZERO
MOV BX, 1
JMP T3_END
T3_NEG:
MOV BX, -1
JMP T3_END
T3_ZERO:
MOV BX, 0
T3_END:
\end{lstlisting}
\subsection*{Task 04 -- If AL is 1 or 3 print o; if 2 or 4 print e}
\begin{lstlisting}[style=asm]
CMP AL,1
JE T4_ODD
CMP AL,3
JE T4_ODD
CMP AL,2
JE T4_EVEN
CMP AL,4
JE T4_EVEN
JMP T4_END
T4_ODD:
MOV DL,'o'
JMP T4_PRINT
T4_EVEN:
MOV DL,'e'
T4_PRINT:
MOV AH,2
INT 21H
T4_END:
\end{lstlisting}
\subsection*{Task 05 -- Read and display only an uppercase letter}
\begin{lstlisting}[style=asm]
MOV AH,1
INT 21H
CMP AL,'A'
JB T5_END
CMP AL,'Z'
JA T5_END
MOV DL,AL
MOV AH,2
INT 21H
T5_END:
\end{lstlisting}
\subsection*{Task 06 -- Display input only if y or Y; otherwise terminate}
\begin{lstlisting}[style=asm]
MOV AH,1
INT 21H
CMP AL,'y'
JE T6_SHOW
CMP AL,'Y'
JNE T6_END
T6_SHOW:
MOV DL,AL
MOV AH,2
INT 21H
T6_END:
\end{lstlisting}
\subsection*{Task 07 -- Check even or odd}
This example checks a byte value stored in AL. Use \texttt{TEST} so the value itself is not changed.
\begin{lstlisting}[style=asm]
MOV AL, 23
TEST AL, 1
JZ T7_EVEN
LEA DX, MSG_ODD
JMP T7_PRINT
T7_EVEN:
LEA DX, MSG_EVEN
T7_PRINT:
MOV AH,9
INT 21H
\end{lstlisting}
Data: \texttt{MSG\_ODD DB 'ODD\$'} and \texttt{MSG\_EVEN DB 'EVEN\$'}.

\subsection*{Task 08 -- Vowel or consonant (case-insensitive)}
\begin{lstlisting}[style=asm]
MOV AH,1
INT 21H
OR AL,20H            ; uppercase ASCII -> lowercase
CMP AL,'a'
JE T8_VOWEL
CMP AL,'e'
JE T8_VOWEL
CMP AL,'i'
JE T8_VOWEL
CMP AL,'o'
JE T8_VOWEL
CMP AL,'u'
JE T8_VOWEL
LEA DX,MSG_CONSONANT
JMP T8_PRINT
T8_VOWEL:
LEA DX,MSG_VOWEL
T8_PRINT:
MOV AH,9
INT 21H
\end{lstlisting}
Data: \texttt{MSG\_VOWEL DB 'VOWEL\$'} and \texttt{MSG\_CONSONANT DB 'CONSONANT\$'}. This assumes the input is an alphabetic character.

\subsection*{Task 09 -- Check divisibility by both 5 and 11}
\begin{lstlisting}[style=asm]
.DATA
NUM DW 55
MSG_Y DB 'Divisible by both 5 and 11$'
MSG_N DB 'Not divisible by both$'
.CODE
MOV AX,NUM
XOR DX,DX
MOV BX,5
DIV BX
CMP DX,0
JNE T9_NO
MOV AX,NUM
XOR DX,DX
MOV BX,11
DIV BX
CMP DX,0
JNE T9_NO
LEA DX,MSG_Y
JMP T9_PRINT
T9_NO:
LEA DX,MSG_N
T9_PRINT:
MOV AH,9
INT 21H
\end{lstlisting}

\subsection*{Task 10 -- Maximum and minimum of three numbers}
The sheet's sample says input 2, 3, 4 but shows minimum 1. That sample minimum is inconsistent: the minimum of 2, 3, 4 is 2. The following correct program uses unsigned word values.
\begin{lstlisting}[style=asm]
.DATA
A DW 2
B DW 3
C DW 4
MAXV DW ?
MINV DW ?
.CODE
MOV AX,A
MOV MAXV,AX
MOV MINV,AX
MOV AX,B
CMP AX,MAXV
JBE T10_CHECK_B_MIN
MOV MAXV,AX
T10_CHECK_B_MIN:
CMP AX,MINV
JAE T10_CHECK_C
MOV MINV,AX
T10_CHECK_C:
MOV AX,C
CMP AX,MAXV
JBE T10_CHECK_C_MIN
MOV MAXV,AX
T10_CHECK_C_MIN:
CMP AX,MINV
JAE T10_DONE
MOV MINV,AX
T10_DONE:
\end{lstlisting}

\subsection*{Task 11 -- Check whether three sides form a valid triangle}
For positive side lengths A, B, C, all three triangle inequalities must hold: A+B>C, A+C>B, B+C>A. This code also handles a 16-bit addition carry as definitely greater than the third unsigned side.
\begin{lstlisting}[style=asm]
.DATA
A DW 3
B DW 4
C DW 5
YESMSG DB 'Y$'
NOMSG DB 'N$'
.CODE
MOV AX,A
ADD AX,B
JC T11_AB_OK
CMP AX,C
JBE T11_NO
T11_AB_OK:
MOV AX,A
ADD AX,C
JC T11_AC_OK
CMP AX,B
JBE T11_NO
T11_AC_OK:
MOV AX,B
ADD AX,C
JC T11_BC_OK
CMP AX,A
JBE T11_NO
T11_BC_OK:
LEA DX,YESMSG
JMP T11_PRINT
T11_NO:
LEA DX,NOMSG
T11_PRINT:
MOV AH,9
INT 21H
\end{lstlisting}
This checks the inequalities but assumes A, B, and C are positive. If reading input, reject zero first.

\subsection*{Task 12 -- Map values 0--3 to i, 4--6 to k, 7--9 to l, and 10 to m}
The sheet says ``a digit'' but includes 10, which cannot be represented by one decimal digit. Treat AX as an integer in 0..10 (or read a two-character decimal number for 10).
\begin{lstlisting}[style=asm]
; AX contains numeric value 0..10
CMP AX,3
JBE T12_I
CMP AX,6
JBE T12_K
CMP AX,9
JBE T12_L
CMP AX,10
JE T12_M
JMP T12_END
T12_I: MOV DL,'i'
       JMP T12_PRINT
T12_K: MOV DL,'k'
       JMP T12_PRINT
T12_L: MOV DL,'l'
       JMP T12_PRINT
T12_M: MOV DL,'m'
T12_PRINT:
MOV AH,2
INT 21H
T12_END:
\end{lstlisting}

\subsection*{Task 13 -- Print weekday; Saturday is day 1}
\begin{lstlisting}[style=asm]
.DATA
D1 DB 'Saturday$'
D2 DB 'Sunday$'
D3 DB 'Monday$'
D4 DB 'Tuesday$'
D5 DB 'Wednesday$'
D6 DB 'Thursday$'
D7 DB 'Friday$'
DAY DB 3
.CODE
MOV AL,DAY
CMP AL,1
JE T13_1
CMP AL,2
JE T13_2
CMP AL,3
JE T13_3
CMP AL,4
JE T13_4
CMP AL,5
JE T13_5
CMP AL,6
JE T13_6
CMP AL,7
JE T13_7
JMP T13_END
T13_1: LEA DX,D1
       JMP T13_PRINT
T13_2: LEA DX,D2
       JMP T13_PRINT
T13_3: LEA DX,D3
       JMP T13_PRINT
T13_4: LEA DX,D4
       JMP T13_PRINT
T13_5: LEA DX,D5
       JMP T13_PRINT
T13_6: LEA DX,D6
       JMP T13_PRINT
T13_7: LEA DX,D7
T13_PRINT:
MOV AH,9
INT 21H
T13_END:
\end{lstlisting}

\subsection*{Task 14 -- Print number of days in a month}
February is treated as 28 days because the task does not provide a year/leap-year input.
\begin{lstlisting}[style=asm]
.DATA
MONTH DB 3
D31 DB '31 day$'
D30 DB '30 day$'
D28 DB '28 day$'
.CODE
MOV AL,MONTH
CMP AL,1
JE T14_31
CMP AL,3
JE T14_31
CMP AL,5
JE T14_31
CMP AL,7
JE T14_31
CMP AL,8
JE T14_31
CMP AL,10
JE T14_31
CMP AL,12
JE T14_31
CMP AL,4
JE T14_30
CMP AL,6
JE T14_30
CMP AL,9
JE T14_30
CMP AL,11
JE T14_30
CMP AL,2
JE T14_28
JMP T14_END
T14_31: LEA DX,D31
        JMP T14_PRINT
T14_30: LEA DX,D30
        JMP T14_PRINT
T14_28: LEA DX,D28
T14_PRINT:
MOV AH,9
INT 21H
T14_END:
\end{lstlisting}

\section{Lab 03(B) -- Looping structures}
\subsection*{Task 01 -- Implicit loop to display 80 stars}
\begin{lstlisting}[style=asm]
MOV CX,80
MOV DL,'*'
T1_LOOP:
    MOV AH,2
    INT 21H
    LOOP T1_LOOP
\end{lstlisting}

\subsection*{Task 02(a) -- Sum 1 + 4 + 7 + ... + 148 into AX}
There are 50 terms. The final value is 3725 decimal (0E8DH).
\begin{lstlisting}[style=asm]
MOV AX,0             ; sum
MOV BX,1             ; current term
MOV CX,50
T2A_LOOP:
    ADD AX,BX
    ADD BX,3
    LOOP T2A_LOOP
; AX = 3725 = 0E8DH
\end{lstlisting}
\subsection*{Task 02(b) -- Sum 100 + 95 + ... + 5 into variable SUM}
There are 20 terms. The final value is 1050 decimal (041AH).
\begin{lstlisting}[style=asm]
.DATA
SUM DW 0
.CODE
MOV AX,0
MOV BX,100
MOV CX,20
T2B_LOOP:
    ADD AX,BX
    SUB BX,5
    LOOP T2B_LOOP
MOV SUM,AX
\end{lstlisting}

\subsection*{Task 03 -- Read five password characters and overprint with five Xs}
The sheet specifies five characters and five X characters. AH=01H echoes typed characters; after five inputs, output CR and five Xs.
\begin{lstlisting}[style=asm]
MOV CX,5
T3_READ:
    MOV AH,1
    INT 21H
    LOOP T3_READ
MOV DL,0DH
MOV AH,2
INT 21H
MOV CX,5
MOV DL,'X'
T3_MASK:
    MOV AH,2
    INT 21H
    LOOP T3_MASK
\end{lstlisting}

\subsection*{Task 04 -- Multiply M and N by repeated addition}
Assume M and N are non-negative and N is a positive count. If N=0, the product is 0 and the loop should be skipped.
\begin{lstlisting}[style=asm]
.DATA
M DW 7
N DW 4
PRODUCT DW 0
.CODE
MOV AX,0
MOV CX,N
CMP CX,0
JE T4_DONE
T4_LOOP:
    ADD AX,M
    LOOP T4_LOOP
T4_DONE:
MOV PRODUCT,AX
\end{lstlisting}

\subsection*{Task 05 -- Display extended ASCII codes 80H through FFH, 10 per line}
This displays each code once, with a blank between characters and CR/LF after every ten. Display appearance for 80H--FFH depends on the active DOS code page/EMU8086 character renderer.
\begin{lstlisting}[style=asm]
MOV SI,80H
MOV BL,0             ; number printed on current line
MOV CX,128
T5_LOOP:
    MOV DX,SI         ; DL = current character code
    MOV AH,2
    INT 21H
    MOV DL,' '
    MOV AH,2
    INT 21H
    INC BL
    CMP BL,10
    JNE T5_NEXT
    MOV DL,0DH
    MOV AH,2
    INT 21H
    MOV DL,0AH
    INT 21H
    MOV BL,0
T5_NEXT:
    INC SI
    LOOP T5_LOOP
\end{lstlisting}

\subsection*{Task 06 -- Hex digit converter with validation and repeat prompt}
\begin{lstlisting}[style=asm]
.DATA
P DB 'ENTER A HEX DIGIT (0-9,A-F): $'
O DB 0DH,0AH,'IN DECIMAL IT IS $'
AGAIN DB 0DH,0AH,'DO YOU WANT TO DO IT AGAIN? $'
BAD DB 0DH,0AH,'ILLEGAL CHARACTER. TRY AGAIN.$'
.CODE
T6_INPUT:
    LEA DX,P
    MOV AH,9
    INT 21H
    MOV AH,1
    INT 21H
    CMP AL,'0'
    JB T6_INVALID
    CMP AL,'9'
    JBE T6_NUMERIC
    CMP AL,'A'
    JB T6_INVALID
    CMP AL,'F'
    JA T6_INVALID
    SUB AL,37H          ; A..F -> 10..15
    JMP T6_SHOW
T6_NUMERIC:
    SUB AL,'0'
T6_SHOW:
    XOR AH,AH
    MOV BX,AX           ; preserve numeric result
    LEA DX,O
    MOV AH,9
    INT 21H
    MOV AX,BX
    CALL PRINT_U16
    LEA DX,AGAIN
    MOV AH,9
    INT 21H
    MOV AH,1
    INT 21H
    CMP AL,'y'
    JE T6_INPUT
    CMP AL,'Y'
    JE T6_INPUT
    JMP T6_END
T6_INVALID:
    LEA DX,BAD
    MOV AH,9
    INT 21H
    JMP T6_INPUT
T6_END:
\end{lstlisting}

\subsection*{Task 07 -- Longest consecutive alphabetically increasing capital-letter run}
Example: \texttt{FGHADEFGHC} produces \texttt{DEFGH}. The program reads until CR and keeps the first longest run when there is a tie.
\begin{lstlisting}[style=asm]
.DATA
PREV DB ?
CURR DB ?
CURSTART DB ?
BESTSTART DB ?
CURLEN DB 0
BESTLEN DB 0
HAVE DB 0
MSG DB 0DH,0AH,'THE LONGEST CONSECUTIVELY INCREASING STRING IS: $'
.CODE
T7_READ:
    MOV AH,1
    INT 21H
    CMP AL,0DH
    JE T7_DONE
    MOV CURR,AL
    CMP HAVE,0
    JNE T7_COMPARE
    MOV HAVE,1
    MOV AL,CURR
    MOV PREV,AL
    MOV CURSTART,AL
    MOV CURLEN,1
    MOV BESTSTART,AL
    MOV BESTLEN,1
    JMP T7_READ
T7_COMPARE:
    MOV BL,PREV
    INC BL
    MOV AL,CURR
    CMP AL,BL
    JE T7_CONTINUE
    MOV AL,CURR
    MOV CURSTART,AL
    MOV CURLEN,1
    JMP T7_UPDATE
T7_CONTINUE:
    INC CURLEN
T7_UPDATE:
    MOV AL,CURLEN
    CMP AL,BESTLEN
    JBE T7_SAVE_PREV
    MOV BESTLEN,AL
    MOV AL,CURSTART
    MOV BESTSTART,AL
T7_SAVE_PREV:
    MOV AL,CURR
    MOV PREV,AL
    JMP T7_READ
T7_DONE:
    LEA DX,MSG
    MOV AH,9
    INT 21H
    MOV DL,BESTSTART
    MOV CL,BESTLEN
    CMP CL,0
    JE T7_END
T7_PRINT:
    MOV AH,2
    INT 21H
    INC DL
    DEC CL
    JNZ T7_PRINT
T7_END:
\end{lstlisting}

\subsection*{Task 08 -- Sum 1..100 divisible by the last non-zero ID digit into M, others into N}
Set KEY to the last non-zero digit of the student's ID (1--9). M is the divisible sum; N is the non-divisible sum.
\begin{lstlisting}[style=asm]
.DATA
KEY DB 3
M DW 0
N DW 0
.CODE
MOV SI,1
MOV M,0
MOV N,0
T8_LOOP:
    MOV AX,SI
    XOR AH,AH          ; SI is 1..100; AX contains the value
    DIV KEY            ; byte divisor; remainder in AH
    CMP AH,0
    JNE T8_OTHER
    MOV AX,M
    ADD AX,SI
    MOV M,AX
    JMP T8_NEXT
T8_OTHER:
    MOV AX,N
    ADD AX,SI
    MOV N,AX
T8_NEXT:
    INC SI
    CMP SI,101
    JB T8_LOOP
\end{lstlisting}

\subsection*{Task 09 -- Nested-loop number pattern}
For input row count 1--9, print row 1 as 1, row 2 as 12, ..., row N as 123...N. This version reads a single input digit.
\begin{lstlisting}[style=asm]
MOV AH,1
INT 21H
SUB AL,'0'
MOV BL,AL            ; number of rows
MOV BH,1              ; current row
T9_OUTER:
    MOV CL,1          ; current column
T9_INNER:
    MOV DL,CL
    ADD DL,'0'
    MOV AH,2
    INT 21H
    INC CL
    CMP CL,BH
    JBE T9_INNER
    MOV DL,0DH
    MOV AH,2
    INT 21H
    MOV DL,0AH
    INT 21H
    INC BH
    CMP BH,BL
    JBE T9_OUTER
\end{lstlisting}

\subsection*{Task 10 -- Read four year digits using a loop and test leap year}
The lab sheet's stated rule has a logical typo: ``divisible by 400 but not by 100'' is impossible, because every multiple of 400 is also a multiple of 100. The correct Gregorian rule is: a year is leap if it is divisible by 400, OR if it is divisible by 4 and not divisible by 100. This complete program body implements that rule.
\begin{lstlisting}[style=asm]
.DATA
YEAR DW 0
DIGIT DB ?
MSG_LEAP DB 0DH,0AH,'Leap year$'
MSG_NOT DB 0DH,0AH,'Not leap year$'
.CODE
MOV CX,4
MOV YEAR,0
T10_READ:
    MOV AH,1
    INT 21H
    SUB AL,'0'
    MOV DIGIT,AL
    MOV AX,YEAR
    MOV BX,10
    MUL BX              ; AX = current year * 10
    XOR BX,BX
    MOV BL,DIGIT
    ADD AX,BX
    MOV YEAR,AX
    LOOP T10_READ

; if divisible by 400 -> leap
MOV AX,YEAR
XOR DX,DX
MOV BX,400
DIV BX
CMP DX,0
JE T10_LEAP
; if not divisible by 4 -> not leap
MOV AX,YEAR
XOR DX,DX
MOV BX,4
DIV BX
CMP DX,0
JNE T10_NOT
; divisible by 4: if divisible by 100 -> not leap
MOV AX,YEAR
XOR DX,DX
MOV BX,100
DIV BX
CMP DX,0
JE T10_NOT
T10_LEAP:
LEA DX,MSG_LEAP
JMP T10_PRINT
T10_NOT:
LEA DX,MSG_NOT
T10_PRINT:
MOV AH,9
INT 21H
\end{lstlisting}
\section{Sample Midterm Question Paper 1 -- Complete worked answers}
\subsection*{Question 1 -- Evaluate ((P+Q) // (R mod S)) * T}
Use variables only. P=14, Q=45, R=9, S=second non-zero student-ID digit, T=last non-zero student-ID digit. Store result in RESULT.
\begin{lstlisting}[style=asm]
.DATA
P DW 14
Q DW 45
R DW 9
S DW 2          ; replace with the required ID digit
T DW 3          ; replace with the required ID digit
RESULT DW ?
ERR DB 'R mod S is zero; expression undefined.$'
.CODE
; remainder = R mod S
MOV AX,R
XOR DX,DX
MOV BX,S
CMP BX,0
JE SQ1_ERROR
DIV BX
MOV BX,DX          ; BX = R mod S
CMP BX,0
JE SQ1_ERROR
; quotient = (P+Q) / (R mod S)
MOV AX,P
ADD AX,Q
XOR DX,DX
DIV BX
; multiply quotient by T
MUL T
MOV RESULT,AX      ; assuming product fits in 16 bits
JMP SQ1_DONE
SQ1_ERROR:
LEA DX,ERR
MOV AH,9
INT 21H
SQ1_DONE:
\end{lstlisting}
If S divides 9 exactly, then R mod S = 0 and the expression divides by zero. In that case there is no defined numeric answer. If the product exceeds 65535, store the full 32-bit result as DX:AX / two words.

\subsection*{Question 2 -- Read and evaluate five signed one-digit expressions}
Input format per line: [sign][digit][operator][sign][digit]. Operators: +, -, *, /. The full program below repeats five times; signed division truncates toward zero. Division by zero prints an error and continues. Add the listed data fields and the procedures after MAIN.
\begin{lstlisting}[style=asm]
.DATA
COUNT DB 5
SIGN1 DB ?
SIGN2 DB ?
OPR DB ?
DIGIT DB ?
OPA DW ?
OPB DW ?
RES DW ?
ERRDIV DB '  DIVISION BY ZERO$'
.CODE
MOV CL,5
SQ2_NEXT:
    ; first sign
    MOV AH,1
    INT 21H
    MOV SIGN1,AL
    ; first digit
    MOV AH,1
    INT 21H
    SUB AL,'0'
    XOR AH,AH
    MOV BX,AX
    CMP SIGN1,'-'
    JNE SQ2_SAVE_A
    NEG BX
SQ2_SAVE_A:
    MOV OPA,BX
    ; operator
    MOV AH,1
    INT 21H
    MOV OPR,AL
    ; second sign
    MOV AH,1
    INT 21H
    MOV SIGN2,AL
    ; second digit
    MOV AH,1
    INT 21H
    SUB AL,'0'
    XOR AH,AH
    MOV BX,AX
    CMP SIGN2,'-'
    JNE SQ2_SAVE_B
    NEG BX
SQ2_SAVE_B:
    MOV OPB,BX
    ; evaluate
    MOV AX,OPA
    CMP OPR,'+'
    JE SQ2_ADD
    CMP OPR,'-'
    JE SQ2_SUB
    CMP OPR,'*'
    JE SQ2_MUL
    CMP OPR,'/'
    JE SQ2_DIV
    JMP SQ2_DISPLAY
SQ2_ADD:
    ADD AX,OPB
    JMP SQ2_STORE
SQ2_SUB:
    SUB AX,OPB
    JMP SQ2_STORE
SQ2_MUL:
    IMUL OPB
    JMP SQ2_STORE
SQ2_DIV:
    CMP OPB,0
    JE SQ2_DIVZERO
    CWD
    IDIV OPB
    JMP SQ2_STORE
SQ2_DIVZERO:
    LEA DX,ERRDIV
    MOV AH,9
    INT 21H
    JMP SQ2_LINE_END
SQ2_STORE:
    MOV RES,AX
SQ2_DISPLAY:
    MOV AX,RES
    CALL PRINT_S16
SQ2_LINE_END:
    CALL NEWLINE
    DEC CL
    JZ SQ2_FINISH
    ; consume the Enter key after the 5-character expression without echo
    MOV AH,8
    INT 21H
    JMP SQ2_NEXT
SQ2_FINISH:
\end{lstlisting}
This assumes the user presses Enter after each five-character expression. If the assessment uses a continuous 25-character stream, remove the CR-consumption call. Include \texttt{PRINT\_S16}, \texttt{PRINT\_U16}, and \texttt{NEWLINE} from Section 1. The sample \texttt{-5++4} evaluates to -1.

\subsection*{Question 3 -- Three command characters control the next three characters}
\begin{lstlisting}[style=asm]
.DATA
C1 DB ?
C2 DB ?
C3 DB ?
X1 DB ?
X2 DB ?
X3 DB ?
.CODE
MOV AH,1
INT 21H
MOV C1,AL
MOV AH,1
INT 21H
MOV C2,AL
MOV AH,1
INT 21H
MOV C3,AL
MOV AH,1
INT 21H
MOV X1,AL
MOV AH,1
INT 21H
MOV X2,AL
MOV AH,1
INT 21H
MOV X3,AL

; command 1 controls character 1
MOV AL,X1
CMP C1,'U'
JNE SQ3_C1_LOWER
CMP AL,'a'
JB SQ3_SAVE1
CMP AL,'z'
JA SQ3_SAVE1
SUB AL,20H
JMP SQ3_SAVE1
SQ3_C1_LOWER:
CMP C1,'L'
JNE SQ3_SAVE1
CMP AL,'A'
JB SQ3_SAVE1
CMP AL,'Z'
JA SQ3_SAVE1
ADD AL,20H
SQ3_SAVE1:
MOV X1,AL

; command 2 controls character 2
MOV AL,X2
CMP C2,'U'
JNE SQ3_C2_LOWER
CMP AL,'a'
JB SQ3_SAVE2
CMP AL,'z'
JA SQ3_SAVE2
SUB AL,20H
JMP SQ3_SAVE2
SQ3_C2_LOWER:
CMP C2,'L'
JNE SQ3_SAVE2
CMP AL,'A'
JB SQ3_SAVE2
CMP AL,'Z'
JA SQ3_SAVE2
ADD AL,20H
SQ3_SAVE2:
MOV X2,AL

; command 3 controls character 3
MOV AL,X3
CMP C3,'U'
JNE SQ3_C3_LOWER
CMP AL,'a'
JB SQ3_SAVE3
CMP AL,'z'
JA SQ3_SAVE3
SUB AL,20H
JMP SQ3_SAVE3
SQ3_C3_LOWER:
CMP C3,'L'
JNE SQ3_SAVE3
CMP AL,'A'
JB SQ3_SAVE3
CMP AL,'Z'
JA SQ3_SAVE3
ADD AL,20H
SQ3_SAVE3:
MOV X3,AL

MOV DL,X1
MOV AH,2
INT 21H
MOV DL,X2
INT 21H
MOV DL,X3
INT 21H
\end{lstlisting}
For input \texttt{UULaBc}, output is \texttt{ABc}. A command other than U or L leaves its character unchanged.

\subsection*{Question 4 -- Sum 1..100 into M (not divisible by key) and N (divisible)}
\begin{lstlisting}[style=asm]
.DATA
PROMPT DB 'Enter the selection key (1-9): $'
M DW 0                  ; not divisible
N DW 0                  ; divisible
KEY DB ?
LABM DB 0DH,0AH,'Starting XI score: $'
LABN DB 0DH,0AH,'Substitute score: $'
.CODE
LEA DX,PROMPT
MOV AH,9
INT 21H
MOV AH,1
INT 21H
SUB AL,'0'
MOV KEY,AL
MOV M,0
MOV N,0
MOV SI,1
SQ4_LOOP:
    MOV AX,SI
    DIV KEY             ; SI is 1..100, AL quotient, AH remainder
    CMP AH,0
    JNE SQ4_NOT_DIV
    MOV AX,N
    ADD AX,SI
    MOV N,AX
    JMP SQ4_NEXT
SQ4_NOT_DIV:
    MOV AX,M
    ADD AX,SI
    MOV M,AX
SQ4_NEXT:
    INC SI
    CMP SI,101
    JB SQ4_LOOP
LEA DX,LABM
MOV AH,9
INT 21H
MOV AX,M
CALL PRINT_U16
LEA DX,LABN
MOV AH,9
INT 21H
MOV AX,N
CALL PRINT_U16
\end{lstlisting}
Expected examples: key 3 gives M=3367 and N=1683; key 9 gives M=4456 and N=594. Input key must be 1--9.

\section{Sample Midterm Question Paper 2 -- Complete worked answers}
\subsection*{Question 1 -- Electricity bill expression}
Equation: $((UNITS^2)\div RATE)+(FIXED\_CHARGE\times MONTHS)$. Store the result in RESULT. All operands are variables.
\begin{lstlisting}[style=asm]
.DATA
UNITS DW 12
RATE DW 5
FIXED_CHARGE DW 4    ; replace with 3rd non-zero ID digit
MONTHS DW 3          ; replace with last non-zero ID digit
RESULT DW ?
.CODE
MOV AX,UNITS
MUL UNITS            ; DX:AX = UNITS squared
MOV BX,RATE
DIV BX               ; AX = floor(UNITS^2 / RATE)
MOV CX,AX             ; save first term
MOV AX,FIXED_CHARGE
MUL MONTHS             ; DX:AX = FIXED_CHARGE * MONTHS
ADD AX,CX
MOV RESULT,AX         ; assuming final value fits in 16 bits
\end{lstlisting}
With the sample values FIXED\_CHARGE=4 and MONTHS=3: $144/5=28$ (integer division), $4\times3=12$, total RESULT=40.

\subsection*{Question 2 -- Convert 24-hour AX to a two-digit 12-hour display}
The input hour is initially in AX (0--23). Output order is hour first, followed by AM/PM, with a leading zero for 1--9.
\begin{lstlisting}[style=asm]
.DATA
AMMSG DB ' AM$'
PMMSG DB ' PM$'
HOUR12 DB ?
ISPM DB ?
.CODE
; AX contains input hour 0..23
CMP AX,0
JNE S2Q2_NOT_ZERO
MOV HOUR12,12
MOV ISPM,0
JMP S2Q2_PRINT_HOUR
S2Q2_NOT_ZERO:
CMP AX,12
JE S2Q2_NOON
JB S2Q2_AM
SUB AX,12
MOV HOUR12,AL
MOV ISPM,1
JMP S2Q2_PRINT_HOUR
S2Q2_NOON:
MOV HOUR12,12
MOV ISPM,1
JMP S2Q2_PRINT_HOUR
S2Q2_AM:
MOV HOUR12,AL
MOV ISPM,0
S2Q2_PRINT_HOUR:
MOV AL,HOUR12
XOR AH,AH
MOV BL,10
DIV BL               ; AL=tens, AH=ones
MOV BH,AH
MOV DL,AL
ADD DL,'0'
MOV AH,2
INT 21H
MOV DL,BH
ADD DL,'0'
MOV AH,2
INT 21H
CMP ISPM,0
JNE S2Q2_PM
LEA DX,AMMSG
JMP S2Q2_SUFFIX
S2Q2_PM:
LEA DX,PMMSG
S2Q2_SUFFIX:
MOV AH,9
INT 21H
\end{lstlisting}
Expected displays: AX=0 -> 12 AM; AX=9 -> 09 AM; AX=12 -> 12 PM; AX=23 -> 11 PM.

\subsection*{Question 3 -- Find second smallest of four register values and leave it in DX}
Use an ascending sorting network with five compare-and-swap pairs. This preserves all four values until the final assignment. The jumps are unsigned because the input range is 1..65535.
\begin{lstlisting}[style=asm]
; Initial values are already in AX, BX, CX, DX.
; Compare-swap AX,BX
CMP AX,BX
JBE S2Q3_AB
XCHG AX,BX
S2Q3_AB:
; Compare-swap CX,DX
CMP CX,DX
JBE S2Q3_CD
XCHG CX,DX
S2Q3_CD:
; Compare-swap AX,CX
CMP AX,CX
JBE S2Q3_AC
XCHG AX,CX
S2Q3_AC:
; Compare-swap BX,DX
CMP BX,DX
JBE S2Q3_BD
XCHG BX,DX
S2Q3_BD:
; Compare-swap BX,CX
CMP BX,CX
JBE S2Q3_BC
XCHG BX,CX
S2Q3_BC:
; Sorted: AX <= BX <= CX <= DX
MOV DX,BX             ; second smallest
\end{lstlisting}
Test 1,23,456,7890 -> DX=23. Test 12345,6789,10,6789 -> DX=6789. Duplicate values are handled correctly.

\subsection*{Question 4 -- Sum squares above and below/equal threshold}
Threshold key is read as one digit. M is sum of squares for numbers greater than threshold; N is sum of squares for numbers less than or equal to threshold.
\begin{lstlisting}[style=asm]
.DATA
PROMPT DB 'Enter the threshold key: $'
ABOVE_MSG DB 0DH,0AH,'Above threshold sum: $'
BELOW_MSG DB 0DH,0AH,'Below threshold sum: $'
KEY DB ?
I DW 1
M DW 0
N DW 0
.CODE
LEA DX,PROMPT
MOV AH,9
INT 21H
MOV AH,1
INT 21H
SUB AL,'0'
MOV KEY,AL
MOV I,1
MOV M,0
MOV N,0
S2Q4_LOOP:
    MOV AX,I
    MUL I                 ; DX:AX = I*I; here I<=50, so fits in AX
    MOV BX,I
    XOR BH,BH             ; BX = I
    CMP BL,KEY
    JBE S2Q4_BELOW
    ADD M,AX
    JMP S2Q4_NEXT
S2Q4_BELOW:
    ADD N,AX
S2Q4_NEXT:
    INC I
    CMP I,51
    JB S2Q4_LOOP
LEA DX,ABOVE_MSG
MOV AH,9
INT 21H
MOV AX,M
CALL PRINT_U16
LEA DX,BELOW_MSG
MOV AH,9
INT 21H
MOV AX,N
CALL PRINT_U16
\end{lstlisting}
For threshold 5, M=42870 and N=55. The total sum of squares $1^2$ through $50^2$ is 42925. The sample sheet prints 41845 for M, but this is an arithmetic error: 42925-55=42870.

\section{Important corrections and assumptions in the supplied materials}
\begin{enumerate}
\item \textbf{CR/LF labels:} correct values are CR=0DH (13 decimal) and LF=0AH (10 decimal). Some extracted Lab 02 text transposes their decimal labels.
\item \textbf{AF for 80H+80H:} AF=0, not 1, because the low nibble is 0+0. Other flags: CF=1, PF=1, ZF=1, SF=0, OF=1.
\item \textbf{Lab 03(B) leap year wording:} ``divisible by 400 but not by 100'' cannot occur. Use the Gregorian rule: divisible by 400 OR divisible by 4 and not by 100.
\item \textbf{Lab 03(A) Task 10 sample:} input 2,3,4 has minimum 2, not the printed 1.
\item \textbf{Lab 03(A) Task 12:} a single decimal digit cannot equal 10. Treat input as an integer 0..10 or use two characters for 10.
\item \textbf{Sample Midterm 2 Question 4:} for threshold 5, the correct sum of squares above the threshold is 42870; the paper prints 41845. The below/equal sum 55 is correct.
\item \textbf{Sample Midterm 1 Question 1:} if S divides R exactly, R mod S is zero and the expression is undefined due to division by zero.
\end{enumerate}

\section{Exam practice checklist}
\begin{itemize}
\item[\(\square\)] I can write the supplied template and initialize DS without looking.
\item[\(\square\)] I know the registers used implicitly by MUL and DIV.
\item[\(\square\)] I can convert a keyboard digit from ASCII to numeric and back.
\item[\(\square\)] I know INT 21H functions 01H, 02H, 09H, and 4CH.
\item[\(\square\)] I can calculate CF, PF, AF, ZF, SF, and OF for a short operation.
\item[\(\square\)] I can choose signed or unsigned jumps correctly.
\item[\(\square\)] I can write IF, IF-ELSE, CASE, explicit, LOOP/CX, repeat-until, and nested-loop structures.
\item[\(\square\)] I can solve a sum/counting task and display a multi-digit decimal result.
\item[\(\square\)] I have traced each solution with Single Step at least once.
\end{itemize}
\begin{center}\Large\bfseries Do not memorize only the code. Understand what each register contains at every label.\end{center}
\end{document}
